%% ---------------------------------------------------------------------------
%% Shared preamble for every standalone figure in this directory.
%%
%% Usage, from figures/experiment/ or figures/schematic/:
%%
%%     \documentclass[crop,tikz,10pt]{standalone}
%%     %% ---------------------------------------------------------------------------
%% Shared preamble for every standalone figure in this directory.
%%
%% Usage, from figures/experiment/ or figures/schematic/:
%%
%%     \documentclass[crop,tikz,10pt]{standalone}
%%     %% ---------------------------------------------------------------------------
%% Shared preamble for every standalone figure in this directory.
%%
%% Usage, from figures/experiment/ or figures/schematic/:
%%
%%     \documentclass[crop,tikz,10pt]{standalone}
%%     %% ---------------------------------------------------------------------------
%% Shared preamble for every standalone figure in this directory.
%%
%% Usage, from figures/experiment/ or figures/schematic/:
%%
%%     \documentclass[crop,tikz,10pt]{standalone}
%%     \input{../preamble}
%%     \begin{document}
%%     \begin{tikzpicture} ... \end{tikzpicture}
%%     \end{document}
%%
%% The whole point of this file is that a figure compiled on its own comes out
%% metrically identical to the same figure compiled inside main.tex.  That
%% requires reproducing three things acmart would otherwise have provided:
%% the font family, the font *sizes*, and \columnwidth / \textwidth.
%% ---------------------------------------------------------------------------

%% --- fonts: exactly what acmart.cls loads for sigconf (see acmart.cls:802) ---
\usepackage[T1]{fontenc}
\usepackage[tt=false, type1=true]{libertine}
\usepackage[varqu]{zi4}
\usepackage[libertine]{newtxmath}

%% --- font sizes: acmart's 10pt values, which are NOT article's 10pt values ---
%% Measured from a real acmart[sigconf,10pt] document.  article would give
%% \footnotesize 8/9.5 and \tiny 5/6; acmart gives 8/10 and 6/7.  Figures lean
%% heavily on \tiny and \scriptsize, so getting this wrong is immediately
%% visible as figure labels that do not match the body text.
\makeatletter
\renewcommand\normalsize{\@setfontsize\normalsize{10}{12}}
\renewcommand\small{\@setfontsize\small{9}{11}}
\renewcommand\footnotesize{\@setfontsize\footnotesize{8}{10}}
\renewcommand\scriptsize{\@setfontsize\scriptsize{7}{8}}
\renewcommand\tiny{\@setfontsize\tiny{6}{7}}
\makeatother
\normalsize

%% --- page geometry: acmart sigconf 10pt, measured, not guessed ---
\newlength{\figcolwidth}   \setlength{\figcolwidth}{241.14749pt}
\newlength{\figtextwidth}  \setlength{\figtextwidth}{506.295pt}
\setlength{\columnsep}{24pt}

%% \figlinewidth is the width of the container this figure sits in, and it is
%% what every plot must size itself against.
%%
%% It exists because \columnwidth and \linewidth are NOT usable here: standalone
%% resets both at \begin{document} to the cropped page, so a style written as
%% `width=1.0\columnwidth' silently becomes the full text width.  \figlinewidth
%% is ours and nothing else touches it.
%%
%% Default is a plain single-column figure.  A figure that lived inside a
%% subfigure says so instead --- in main.tex a subfigure is a minipage, and
%% LaTeX's minipage sets \linewidth/\columnwidth/\textwidth to the minipage
%% width, so `width=1.0\linewidth' there was measuring the subfigure:
%%
%%     \figcontainer{0.48\figtextwidth}   % subfigure{0.48\textwidth} in figure*
%%     \figcontainer{0.49\figcolwidth}    % subfigure{0.49\columnwidth}
%%
%% Call it in the document body, before the picture.  Making this explicit is
%% strictly better than the inheritance it replaces: the figure now states the
%% width it was designed for instead of inheriting it from a wrapper two files
%% away.
\newlength{\figlinewidth}  \setlength{\figlinewidth}{\figcolwidth}
%% \ignorespaces is load-bearing: without it the newline after a
%% \figcontainer{...} line leaves an interword space in the document body,
%% and standalone's crop counts that ~2.5pt of glue as part of the figure ---
%% the picture comes out wider than its container and lands shifted right.
\newcommand{\figcontainer}[1]{\setlength{\figlinewidth}{#1}\ignorespaces}

\input{../macros}

\usepackage{amsmath}
\usepackage{pgfplots}
\pgfplotsset{compat=1.18}
\usetikzlibrary{positioning,arrows.meta,fit,backgrounds,calc,patterns,automata}

%% ---------------------------------------------------------------------------
%% Plot house style.  Conventions used consistently in every figure:
%%   cB (red, thick, filled circles) .... ActDelta, always
%%   cA (blue, dashed, squares) ......... the arm we are compared against
%%   cGrey .............................. status quo, oracles, annotations
%%   shaded band ........................ a region where something fails
%% Series are separable in grayscale by dash pattern + marker, not colour alone.
%% ---------------------------------------------------------------------------
\pgfplotsset{
  compactplot/.style={
    width=1.0\figlinewidth, height=4.1cm,
    tick label style={font=\scriptsize},
    label style={font=\small},
    legend style={font=\scriptsize, draw=none, fill=white, fill opacity=0.88,
      text opacity=1, inner xsep=2pt, inner ysep=1pt,
      nodes={inner sep=1pt}, cells={anchor=west}},
    title style={font=\scriptsize, yshift=-1.2mm},
    grid=major, grid style={gray!18, line width=0.25pt},
    tick align=outside, tick pos=left,
    axis line style={gray!55, line width=0.5pt},
    every axis plot/.append style={line width=0.9pt},
    legend cell align=left,
  },
  widecompact/.style={
    compactplot, width=1.0\figlinewidth, height=4.3cm,
  },
}

%% "higher/lower is better" arrow.  Carried over from the old main.tex
%% preamble, where it was defined but never called --- the Pareto plot draws
%% its arrow inline at 1pt, and this macro draws at 0.9pt.  Kept so the two do
%% not silently converge; if you adopt it, adopt it everywhere and accept that
%% the Pareto arrow gets thinner.
\newcommand{\betterarrow}[4]{%
  \draw[-{Latex[length=1.8mm]}, cGrey, line width=0.9pt]
    (axis cs:#1,#2) -- (axis cs:#3,#4);}

%% ---------------------------------------------------------------------------
%% Data access.  build.py reads <figure>.json and writes ../.build/<figure>/,
%% which this preamble pulls in automatically via \jobname.  A figure's .tex
%% therefore never contains a measured number: it asks for one by name.
%%
%%   \figdata{series}   -> path to that series' .dat, for \addplot table
%%   \figval{key}       -> a scalar from the JSON "scalars" block
%%   \figlabel{key}     -> a series label from the JSON "labels" block
%%   \figgrid{key}      -> generated \heatcell calls for a heatmap
%%   \figgridcols{key}  -> "0/{0-12},1/{12-25},..." for the x tick \foreach
%%   \figgridrows{key}  -> "0/50,1/100,..."        for the y tick \foreach
%% ---------------------------------------------------------------------------
\newcommand{\figbuild}{../.build/\jobname}
\newcommand{\figdata}[1]{\figbuild/#1.dat}

\makeatletter
\newcommand{\fig@def}[3]{\expandafter\gdef\csname fig@#1@#2\endcsname{#3}}
%% Every line end inside this macro needs its %%: without them the newlines
%% become space tokens, so \figlabel{x} expands to " x " and every legend entry
%% and node label comes out a space-width wider than it was written.
\newcommand{\fig@get}[2]{%
  \ifcsname fig@#1@#2\endcsname%
    \csname fig@#1@#2\endcsname%
  \else%
    \PackageError{figdata}{No #1 named `#2' in \jobname.json}%
      {Add it to the JSON, then rerun `make'.}%
  \fi}
\newcommand{\figvaldef}[2]{\fig@def{val}{#1}{#2}}
\newcommand{\figlabeldef}[2]{\fig@def{label}{#1}{#2}}
\newcommand{\figgriddef}[2]{\fig@def{grid}{#1}{#2}}
\newcommand{\figgridcolsdef}[2]{\fig@def{gridcols}{#1}{#2}}
\newcommand{\figgridrowsdef}[2]{\fig@def{gridrows}{#1}{#2}}
\newcommand{\figval}[1]{\fig@get{val}{#1}}
\newcommand{\figlabel}[1]{\fig@get{label}{#1}}
\newcommand{\figgrid}[1]{\fig@get{grid}{#1}}
\newcommand{\figgridcols}[1]{\fig@get{gridcols}{#1}}
\newcommand{\figgridrows}[1]{\fig@get{gridrows}{#1}}
\makeatother

\InputIfFileExists{\figbuild/data.tex}{}{%
  \PackageError{figdata}{\figbuild/data.tex missing}%
    {Run `make' in figures/ --- build.py turns \jobname.json into it.}}

%%     \begin{document}
%%     \begin{tikzpicture} ... \end{tikzpicture}
%%     \end{document}
%%
%% The whole point of this file is that a figure compiled on its own comes out
%% metrically identical to the same figure compiled inside main.tex.  That
%% requires reproducing three things acmart would otherwise have provided:
%% the font family, the font *sizes*, and \columnwidth / \textwidth.
%% ---------------------------------------------------------------------------

%% --- fonts: exactly what acmart.cls loads for sigconf (see acmart.cls:802) ---
\usepackage[T1]{fontenc}
\usepackage[tt=false, type1=true]{libertine}
\usepackage[varqu]{zi4}
\usepackage[libertine]{newtxmath}

%% --- font sizes: acmart's 10pt values, which are NOT article's 10pt values ---
%% Measured from a real acmart[sigconf,10pt] document.  article would give
%% \footnotesize 8/9.5 and \tiny 5/6; acmart gives 8/10 and 6/7.  Figures lean
%% heavily on \tiny and \scriptsize, so getting this wrong is immediately
%% visible as figure labels that do not match the body text.
\makeatletter
\renewcommand\normalsize{\@setfontsize\normalsize{10}{12}}
\renewcommand\small{\@setfontsize\small{9}{11}}
\renewcommand\footnotesize{\@setfontsize\footnotesize{8}{10}}
\renewcommand\scriptsize{\@setfontsize\scriptsize{7}{8}}
\renewcommand\tiny{\@setfontsize\tiny{6}{7}}
\makeatother
\normalsize

%% --- page geometry: acmart sigconf 10pt, measured, not guessed ---
\newlength{\figcolwidth}   \setlength{\figcolwidth}{241.14749pt}
\newlength{\figtextwidth}  \setlength{\figtextwidth}{506.295pt}
\setlength{\columnsep}{24pt}

%% \figlinewidth is the width of the container this figure sits in, and it is
%% what every plot must size itself against.
%%
%% It exists because \columnwidth and \linewidth are NOT usable here: standalone
%% resets both at \begin{document} to the cropped page, so a style written as
%% `width=1.0\columnwidth' silently becomes the full text width.  \figlinewidth
%% is ours and nothing else touches it.
%%
%% Default is a plain single-column figure.  A figure that lived inside a
%% subfigure says so instead --- in main.tex a subfigure is a minipage, and
%% LaTeX's minipage sets \linewidth/\columnwidth/\textwidth to the minipage
%% width, so `width=1.0\linewidth' there was measuring the subfigure:
%%
%%     \figcontainer{0.48\figtextwidth}   % subfigure{0.48\textwidth} in figure*
%%     \figcontainer{0.49\figcolwidth}    % subfigure{0.49\columnwidth}
%%
%% Call it in the document body, before the picture.  Making this explicit is
%% strictly better than the inheritance it replaces: the figure now states the
%% width it was designed for instead of inheriting it from a wrapper two files
%% away.
\newlength{\figlinewidth}  \setlength{\figlinewidth}{\figcolwidth}
%% \ignorespaces is load-bearing: without it the newline after a
%% \figcontainer{...} line leaves an interword space in the document body,
%% and standalone's crop counts that ~2.5pt of glue as part of the figure ---
%% the picture comes out wider than its container and lands shifted right.
\newcommand{\figcontainer}[1]{\setlength{\figlinewidth}{#1}\ignorespaces}

%% ---------------------------------------------------------------------------
%% Everything main.tex AND the standalone figures both need.
%%
%% \input by main.tex and by figures/preamble.tex, so a macro defined here
%% renders identically in the body text, in a caption, and inside a figure.
%% Never define these anywhere else --- a second definition is how figure text
%% and body text drift apart.
%%
%% The rule for what lives here rather than in preamble.tex: if main.tex refers
%% to it, it belongs here.  \cnum is the reason the colours are here too --- it
%% is used both inside Figure 1 and in that figure's caption, and captions stay
%% in main.tex.
%% ---------------------------------------------------------------------------
\usepackage{xspace}

\newcommand{\sys}{ActDelta\xspace}
\newcommand{\pizero}{$\pi_0$\xspace}
\newcommand{\pizerofive}{$\pi_{0.5}$\xspace}
\newcommand{\pp}{\,pp\xspace}

\definecolor{cA}{RGB}{31,119,180}
\definecolor{cB}{RGB}{214,39,40}
\definecolor{cC}{RGB}{44,160,44}
\definecolor{cD}{RGB}{255,127,14}
\definecolor{cE}{RGB}{148,103,189}
\definecolor{cGrey}{RGB}{120,120,120}

%% circled step number, safe inside captions (no tikz in moving arguments)
\newcommand{\cnum}[1]{\textcolor{cB}{\textcircled{\raisebox{0.35pt}{\tiny #1}}}}


\usepackage{amsmath}
\usepackage{pgfplots}
\pgfplotsset{compat=1.18}
\usetikzlibrary{positioning,arrows.meta,fit,backgrounds,calc,patterns,automata}

%% ---------------------------------------------------------------------------
%% Plot house style.  Conventions used consistently in every figure:
%%   cB (red, thick, filled circles) .... ActDelta, always
%%   cA (blue, dashed, squares) ......... the arm we are compared against
%%   cGrey .............................. status quo, oracles, annotations
%%   shaded band ........................ a region where something fails
%% Series are separable in grayscale by dash pattern + marker, not colour alone.
%% ---------------------------------------------------------------------------
\pgfplotsset{
  compactplot/.style={
    width=1.0\figlinewidth, height=4.1cm,
    tick label style={font=\scriptsize},
    label style={font=\small},
    legend style={font=\scriptsize, draw=none, fill=white, fill opacity=0.88,
      text opacity=1, inner xsep=2pt, inner ysep=1pt,
      nodes={inner sep=1pt}, cells={anchor=west}},
    title style={font=\scriptsize, yshift=-1.2mm},
    grid=major, grid style={gray!18, line width=0.25pt},
    tick align=outside, tick pos=left,
    axis line style={gray!55, line width=0.5pt},
    every axis plot/.append style={line width=0.9pt},
    legend cell align=left,
  },
  widecompact/.style={
    compactplot, width=1.0\figlinewidth, height=4.3cm,
  },
}

%% "higher/lower is better" arrow.  Carried over from the old main.tex
%% preamble, where it was defined but never called --- the Pareto plot draws
%% its arrow inline at 1pt, and this macro draws at 0.9pt.  Kept so the two do
%% not silently converge; if you adopt it, adopt it everywhere and accept that
%% the Pareto arrow gets thinner.
\newcommand{\betterarrow}[4]{%
  \draw[-{Latex[length=1.8mm]}, cGrey, line width=0.9pt]
    (axis cs:#1,#2) -- (axis cs:#3,#4);}

%% ---------------------------------------------------------------------------
%% Data access.  build.py reads <figure>.json and writes ../.build/<figure>/,
%% which this preamble pulls in automatically via \jobname.  A figure's .tex
%% therefore never contains a measured number: it asks for one by name.
%%
%%   \figdata{series}   -> path to that series' .dat, for \addplot table
%%   \figval{key}       -> a scalar from the JSON "scalars" block
%%   \figlabel{key}     -> a series label from the JSON "labels" block
%%   \figgrid{key}      -> generated \heatcell calls for a heatmap
%%   \figgridcols{key}  -> "0/{0-12},1/{12-25},..." for the x tick \foreach
%%   \figgridrows{key}  -> "0/50,1/100,..."        for the y tick \foreach
%% ---------------------------------------------------------------------------
\newcommand{\figbuild}{../.build/\jobname}
\newcommand{\figdata}[1]{\figbuild/#1.dat}

\makeatletter
\newcommand{\fig@def}[3]{\expandafter\gdef\csname fig@#1@#2\endcsname{#3}}
%% Every line end inside this macro needs its %%: without them the newlines
%% become space tokens, so \figlabel{x} expands to " x " and every legend entry
%% and node label comes out a space-width wider than it was written.
\newcommand{\fig@get}[2]{%
  \ifcsname fig@#1@#2\endcsname%
    \csname fig@#1@#2\endcsname%
  \else%
    \PackageError{figdata}{No #1 named `#2' in \jobname.json}%
      {Add it to the JSON, then rerun `make'.}%
  \fi}
\newcommand{\figvaldef}[2]{\fig@def{val}{#1}{#2}}
\newcommand{\figlabeldef}[2]{\fig@def{label}{#1}{#2}}
\newcommand{\figgriddef}[2]{\fig@def{grid}{#1}{#2}}
\newcommand{\figgridcolsdef}[2]{\fig@def{gridcols}{#1}{#2}}
\newcommand{\figgridrowsdef}[2]{\fig@def{gridrows}{#1}{#2}}
\newcommand{\figval}[1]{\fig@get{val}{#1}}
\newcommand{\figlabel}[1]{\fig@get{label}{#1}}
\newcommand{\figgrid}[1]{\fig@get{grid}{#1}}
\newcommand{\figgridcols}[1]{\fig@get{gridcols}{#1}}
\newcommand{\figgridrows}[1]{\fig@get{gridrows}{#1}}
\makeatother

\InputIfFileExists{\figbuild/data.tex}{}{%
  \PackageError{figdata}{\figbuild/data.tex missing}%
    {Run `make' in figures/ --- build.py turns \jobname.json into it.}}

%%     \begin{document}
%%     \begin{tikzpicture} ... \end{tikzpicture}
%%     \end{document}
%%
%% The whole point of this file is that a figure compiled on its own comes out
%% metrically identical to the same figure compiled inside main.tex.  That
%% requires reproducing three things acmart would otherwise have provided:
%% the font family, the font *sizes*, and \columnwidth / \textwidth.
%% ---------------------------------------------------------------------------

%% --- fonts: exactly what acmart.cls loads for sigconf (see acmart.cls:802) ---
\usepackage[T1]{fontenc}
\usepackage[tt=false, type1=true]{libertine}
\usepackage[varqu]{zi4}
\usepackage[libertine]{newtxmath}

%% --- font sizes: acmart's 10pt values, which are NOT article's 10pt values ---
%% Measured from a real acmart[sigconf,10pt] document.  article would give
%% \footnotesize 8/9.5 and \tiny 5/6; acmart gives 8/10 and 6/7.  Figures lean
%% heavily on \tiny and \scriptsize, so getting this wrong is immediately
%% visible as figure labels that do not match the body text.
\makeatletter
\renewcommand\normalsize{\@setfontsize\normalsize{10}{12}}
\renewcommand\small{\@setfontsize\small{9}{11}}
\renewcommand\footnotesize{\@setfontsize\footnotesize{8}{10}}
\renewcommand\scriptsize{\@setfontsize\scriptsize{7}{8}}
\renewcommand\tiny{\@setfontsize\tiny{6}{7}}
\makeatother
\normalsize

%% --- page geometry: acmart sigconf 10pt, measured, not guessed ---
\newlength{\figcolwidth}   \setlength{\figcolwidth}{241.14749pt}
\newlength{\figtextwidth}  \setlength{\figtextwidth}{506.295pt}
\setlength{\columnsep}{24pt}

%% \figlinewidth is the width of the container this figure sits in, and it is
%% what every plot must size itself against.
%%
%% It exists because \columnwidth and \linewidth are NOT usable here: standalone
%% resets both at \begin{document} to the cropped page, so a style written as
%% `width=1.0\columnwidth' silently becomes the full text width.  \figlinewidth
%% is ours and nothing else touches it.
%%
%% Default is a plain single-column figure.  A figure that lived inside a
%% subfigure says so instead --- in main.tex a subfigure is a minipage, and
%% LaTeX's minipage sets \linewidth/\columnwidth/\textwidth to the minipage
%% width, so `width=1.0\linewidth' there was measuring the subfigure:
%%
%%     \figcontainer{0.48\figtextwidth}   % subfigure{0.48\textwidth} in figure*
%%     \figcontainer{0.49\figcolwidth}    % subfigure{0.49\columnwidth}
%%
%% Call it in the document body, before the picture.  Making this explicit is
%% strictly better than the inheritance it replaces: the figure now states the
%% width it was designed for instead of inheriting it from a wrapper two files
%% away.
\newlength{\figlinewidth}  \setlength{\figlinewidth}{\figcolwidth}
%% \ignorespaces is load-bearing: without it the newline after a
%% \figcontainer{...} line leaves an interword space in the document body,
%% and standalone's crop counts that ~2.5pt of glue as part of the figure ---
%% the picture comes out wider than its container and lands shifted right.
\newcommand{\figcontainer}[1]{\setlength{\figlinewidth}{#1}\ignorespaces}

%% ---------------------------------------------------------------------------
%% Everything main.tex AND the standalone figures both need.
%%
%% \input by main.tex and by figures/preamble.tex, so a macro defined here
%% renders identically in the body text, in a caption, and inside a figure.
%% Never define these anywhere else --- a second definition is how figure text
%% and body text drift apart.
%%
%% The rule for what lives here rather than in preamble.tex: if main.tex refers
%% to it, it belongs here.  \cnum is the reason the colours are here too --- it
%% is used both inside Figure 1 and in that figure's caption, and captions stay
%% in main.tex.
%% ---------------------------------------------------------------------------
\usepackage{xspace}

\newcommand{\sys}{ActDelta\xspace}
\newcommand{\pizero}{$\pi_0$\xspace}
\newcommand{\pizerofive}{$\pi_{0.5}$\xspace}
\newcommand{\pp}{\,pp\xspace}

\definecolor{cA}{RGB}{31,119,180}
\definecolor{cB}{RGB}{214,39,40}
\definecolor{cC}{RGB}{44,160,44}
\definecolor{cD}{RGB}{255,127,14}
\definecolor{cE}{RGB}{148,103,189}
\definecolor{cGrey}{RGB}{120,120,120}

%% circled step number, safe inside captions (no tikz in moving arguments)
\newcommand{\cnum}[1]{\textcolor{cB}{\textcircled{\raisebox{0.35pt}{\tiny #1}}}}


\usepackage{amsmath}
\usepackage{pgfplots}
\pgfplotsset{compat=1.18}
\usetikzlibrary{positioning,arrows.meta,fit,backgrounds,calc,patterns,automata}

%% ---------------------------------------------------------------------------
%% Plot house style.  Conventions used consistently in every figure:
%%   cB (red, thick, filled circles) .... ActDelta, always
%%   cA (blue, dashed, squares) ......... the arm we are compared against
%%   cGrey .............................. status quo, oracles, annotations
%%   shaded band ........................ a region where something fails
%% Series are separable in grayscale by dash pattern + marker, not colour alone.
%% ---------------------------------------------------------------------------
\pgfplotsset{
  compactplot/.style={
    width=1.0\figlinewidth, height=4.1cm,
    tick label style={font=\scriptsize},
    label style={font=\small},
    legend style={font=\scriptsize, draw=none, fill=white, fill opacity=0.88,
      text opacity=1, inner xsep=2pt, inner ysep=1pt,
      nodes={inner sep=1pt}, cells={anchor=west}},
    title style={font=\scriptsize, yshift=-1.2mm},
    grid=major, grid style={gray!18, line width=0.25pt},
    tick align=outside, tick pos=left,
    axis line style={gray!55, line width=0.5pt},
    every axis plot/.append style={line width=0.9pt},
    legend cell align=left,
  },
  widecompact/.style={
    compactplot, width=1.0\figlinewidth, height=4.3cm,
  },
}

%% "higher/lower is better" arrow.  Carried over from the old main.tex
%% preamble, where it was defined but never called --- the Pareto plot draws
%% its arrow inline at 1pt, and this macro draws at 0.9pt.  Kept so the two do
%% not silently converge; if you adopt it, adopt it everywhere and accept that
%% the Pareto arrow gets thinner.
\newcommand{\betterarrow}[4]{%
  \draw[-{Latex[length=1.8mm]}, cGrey, line width=0.9pt]
    (axis cs:#1,#2) -- (axis cs:#3,#4);}

%% ---------------------------------------------------------------------------
%% Data access.  build.py reads <figure>.json and writes ../.build/<figure>/,
%% which this preamble pulls in automatically via \jobname.  A figure's .tex
%% therefore never contains a measured number: it asks for one by name.
%%
%%   \figdata{series}   -> path to that series' .dat, for \addplot table
%%   \figval{key}       -> a scalar from the JSON "scalars" block
%%   \figlabel{key}     -> a series label from the JSON "labels" block
%%   \figgrid{key}      -> generated \heatcell calls for a heatmap
%%   \figgridcols{key}  -> "0/{0-12},1/{12-25},..." for the x tick \foreach
%%   \figgridrows{key}  -> "0/50,1/100,..."        for the y tick \foreach
%% ---------------------------------------------------------------------------
\newcommand{\figbuild}{../.build/\jobname}
\newcommand{\figdata}[1]{\figbuild/#1.dat}

\makeatletter
\newcommand{\fig@def}[3]{\expandafter\gdef\csname fig@#1@#2\endcsname{#3}}
%% Every line end inside this macro needs its %%: without them the newlines
%% become space tokens, so \figlabel{x} expands to " x " and every legend entry
%% and node label comes out a space-width wider than it was written.
\newcommand{\fig@get}[2]{%
  \ifcsname fig@#1@#2\endcsname%
    \csname fig@#1@#2\endcsname%
  \else%
    \PackageError{figdata}{No #1 named `#2' in \jobname.json}%
      {Add it to the JSON, then rerun `make'.}%
  \fi}
\newcommand{\figvaldef}[2]{\fig@def{val}{#1}{#2}}
\newcommand{\figlabeldef}[2]{\fig@def{label}{#1}{#2}}
\newcommand{\figgriddef}[2]{\fig@def{grid}{#1}{#2}}
\newcommand{\figgridcolsdef}[2]{\fig@def{gridcols}{#1}{#2}}
\newcommand{\figgridrowsdef}[2]{\fig@def{gridrows}{#1}{#2}}
\newcommand{\figval}[1]{\fig@get{val}{#1}}
\newcommand{\figlabel}[1]{\fig@get{label}{#1}}
\newcommand{\figgrid}[1]{\fig@get{grid}{#1}}
\newcommand{\figgridcols}[1]{\fig@get{gridcols}{#1}}
\newcommand{\figgridrows}[1]{\fig@get{gridrows}{#1}}
\makeatother

\InputIfFileExists{\figbuild/data.tex}{}{%
  \PackageError{figdata}{\figbuild/data.tex missing}%
    {Run `make' in figures/ --- build.py turns \jobname.json into it.}}

%%     \begin{document}
%%     \begin{tikzpicture} ... \end{tikzpicture}
%%     \end{document}
%%
%% The whole point of this file is that a figure compiled on its own comes out
%% metrically identical to the same figure compiled inside main.tex.  That
%% requires reproducing three things acmart would otherwise have provided:
%% the font family, the font *sizes*, and \columnwidth / \textwidth.
%% ---------------------------------------------------------------------------

%% --- fonts: exactly what acmart.cls loads for sigconf (see acmart.cls:802) ---
\usepackage[T1]{fontenc}
\usepackage[tt=false, type1=true]{libertine}
\usepackage[varqu]{zi4}
\usepackage[libertine]{newtxmath}

%% --- font sizes: acmart's 10pt values, which are NOT article's 10pt values ---
%% Measured from a real acmart[sigconf,10pt] document.  article would give
%% \footnotesize 8/9.5 and \tiny 5/6; acmart gives 8/10 and 6/7.  Figures lean
%% heavily on \tiny and \scriptsize, so getting this wrong is immediately
%% visible as figure labels that do not match the body text.
\makeatletter
\renewcommand\normalsize{\@setfontsize\normalsize{10}{12}}
\renewcommand\small{\@setfontsize\small{9}{11}}
\renewcommand\footnotesize{\@setfontsize\footnotesize{8}{10}}
\renewcommand\scriptsize{\@setfontsize\scriptsize{7}{8}}
\renewcommand\tiny{\@setfontsize\tiny{6}{7}}
\makeatother
\normalsize

%% --- page geometry: acmart sigconf 10pt, measured, not guessed ---
\newlength{\figcolwidth}   \setlength{\figcolwidth}{241.14749pt}
\newlength{\figtextwidth}  \setlength{\figtextwidth}{506.295pt}
\setlength{\columnsep}{24pt}

%% \figlinewidth is the width of the container this figure sits in, and it is
%% what every plot must size itself against.
%%
%% It exists because \columnwidth and \linewidth are NOT usable here: standalone
%% resets both at \begin{document} to the cropped page, so a style written as
%% `width=1.0\columnwidth' silently becomes the full text width.  \figlinewidth
%% is ours and nothing else touches it.
%%
%% Default is a plain single-column figure.  A figure that lived inside a
%% subfigure says so instead --- in main.tex a subfigure is a minipage, and
%% LaTeX's minipage sets \linewidth/\columnwidth/\textwidth to the minipage
%% width, so `width=1.0\linewidth' there was measuring the subfigure:
%%
%%     \figcontainer{0.48\figtextwidth}   % subfigure{0.48\textwidth} in figure*
%%     \figcontainer{0.49\figcolwidth}    % subfigure{0.49\columnwidth}
%%
%% Call it in the document body, before the picture.  Making this explicit is
%% strictly better than the inheritance it replaces: the figure now states the
%% width it was designed for instead of inheriting it from a wrapper two files
%% away.
\newlength{\figlinewidth}  \setlength{\figlinewidth}{\figcolwidth}
%% \ignorespaces is load-bearing: without it the newline after a
%% \figcontainer{...} line leaves an interword space in the document body,
%% and standalone's crop counts that ~2.5pt of glue as part of the figure ---
%% the picture comes out wider than its container and lands shifted right.
\newcommand{\figcontainer}[1]{\setlength{\figlinewidth}{#1}\ignorespaces}

%% ---------------------------------------------------------------------------
%% Everything main.tex AND the standalone figures both need.
%%
%% \input by main.tex and by figures/preamble.tex, so a macro defined here
%% renders identically in the body text, in a caption, and inside a figure.
%% Never define these anywhere else --- a second definition is how figure text
%% and body text drift apart.
%%
%% The rule for what lives here rather than in preamble.tex: if main.tex refers
%% to it, it belongs here.  \cnum is the reason the colours are here too --- it
%% is used both inside Figure 1 and in that figure's caption, and captions stay
%% in main.tex.
%% ---------------------------------------------------------------------------
\usepackage{xspace}

\newcommand{\sys}{ActDelta\xspace}
\newcommand{\pizero}{$\pi_0$\xspace}
\newcommand{\pizerofive}{$\pi_{0.5}$\xspace}
\newcommand{\pp}{\,pp\xspace}

\definecolor{cA}{RGB}{31,119,180}
\definecolor{cB}{RGB}{214,39,40}
\definecolor{cC}{RGB}{44,160,44}
\definecolor{cD}{RGB}{255,127,14}
\definecolor{cE}{RGB}{148,103,189}
\definecolor{cGrey}{RGB}{120,120,120}

%% circled step number, safe inside captions (no tikz in moving arguments)
\newcommand{\cnum}[1]{\textcolor{cB}{\textcircled{\raisebox{0.35pt}{\tiny #1}}}}


\usepackage{amsmath}
\usepackage{pgfplots}
\pgfplotsset{compat=1.18}
\usetikzlibrary{positioning,arrows.meta,fit,backgrounds,calc,patterns,automata}

%% ---------------------------------------------------------------------------
%% Plot house style.  Conventions used consistently in every figure:
%%   cB (red, thick, filled circles) .... ActDelta, always
%%   cA (blue, dashed, squares) ......... the arm we are compared against
%%   cGrey .............................. status quo, oracles, annotations
%%   shaded band ........................ a region where something fails
%% Series are separable in grayscale by dash pattern + marker, not colour alone.
%% ---------------------------------------------------------------------------
\pgfplotsset{
  compactplot/.style={
    width=1.0\figlinewidth, height=4.1cm,
    tick label style={font=\scriptsize},
    label style={font=\small},
    legend style={font=\scriptsize, draw=none, fill=white, fill opacity=0.88,
      text opacity=1, inner xsep=2pt, inner ysep=1pt,
      nodes={inner sep=1pt}, cells={anchor=west}},
    title style={font=\scriptsize, yshift=-1.2mm},
    grid=major, grid style={gray!18, line width=0.25pt},
    tick align=outside, tick pos=left,
    axis line style={gray!55, line width=0.5pt},
    every axis plot/.append style={line width=0.9pt},
    legend cell align=left,
  },
  widecompact/.style={
    compactplot, width=1.0\figlinewidth, height=4.3cm,
  },
}

%% "higher/lower is better" arrow.  Carried over from the old main.tex
%% preamble, where it was defined but never called --- the Pareto plot draws
%% its arrow inline at 1pt, and this macro draws at 0.9pt.  Kept so the two do
%% not silently converge; if you adopt it, adopt it everywhere and accept that
%% the Pareto arrow gets thinner.
\newcommand{\betterarrow}[4]{%
  \draw[-{Latex[length=1.8mm]}, cGrey, line width=0.9pt]
    (axis cs:#1,#2) -- (axis cs:#3,#4);}

%% ---------------------------------------------------------------------------
%% Data access.  build.py reads <figure>.json and writes ../.build/<figure>/,
%% which this preamble pulls in automatically via \jobname.  A figure's .tex
%% therefore never contains a measured number: it asks for one by name.
%%
%%   \figdata{series}   -> path to that series' .dat, for \addplot table
%%   \figval{key}       -> a scalar from the JSON "scalars" block
%%   \figlabel{key}     -> a series label from the JSON "labels" block
%%   \figgrid{key}      -> generated \heatcell calls for a heatmap
%%   \figgridcols{key}  -> "0/{0-12},1/{12-25},..." for the x tick \foreach
%%   \figgridrows{key}  -> "0/50,1/100,..."        for the y tick \foreach
%% ---------------------------------------------------------------------------
\newcommand{\figbuild}{../.build/\jobname}
\newcommand{\figdata}[1]{\figbuild/#1.dat}

\makeatletter
\newcommand{\fig@def}[3]{\expandafter\gdef\csname fig@#1@#2\endcsname{#3}}
%% Every line end inside this macro needs its %%: without them the newlines
%% become space tokens, so \figlabel{x} expands to " x " and every legend entry
%% and node label comes out a space-width wider than it was written.
\newcommand{\fig@get}[2]{%
  \ifcsname fig@#1@#2\endcsname%
    \csname fig@#1@#2\endcsname%
  \else%
    \PackageError{figdata}{No #1 named `#2' in \jobname.json}%
      {Add it to the JSON, then rerun `make'.}%
  \fi}
\newcommand{\figvaldef}[2]{\fig@def{val}{#1}{#2}}
\newcommand{\figlabeldef}[2]{\fig@def{label}{#1}{#2}}
\newcommand{\figgriddef}[2]{\fig@def{grid}{#1}{#2}}
\newcommand{\figgridcolsdef}[2]{\fig@def{gridcols}{#1}{#2}}
\newcommand{\figgridrowsdef}[2]{\fig@def{gridrows}{#1}{#2}}
\newcommand{\figval}[1]{\fig@get{val}{#1}}
\newcommand{\figlabel}[1]{\fig@get{label}{#1}}
\newcommand{\figgrid}[1]{\fig@get{grid}{#1}}
\newcommand{\figgridcols}[1]{\fig@get{gridcols}{#1}}
\newcommand{\figgridrows}[1]{\fig@get{gridrows}{#1}}
\makeatother

\InputIfFileExists{\figbuild/data.tex}{}{%
  \PackageError{figdata}{\figbuild/data.tex missing}%
    {Run `make' in figures/ --- build.py turns \jobname.json into it.}}
